\section{第一阶段深讲：硬件机制怎样逼出操作系统}

前面已经介绍了许多硬件部件。本节不再按“CPU、内存、总线、设备”逐项罗列，而是换一个角度：每增加一种硬件能力，系统会遇到什么新问题，操作系统为什么必须出现相应的对象、策略和算法。这样读者不会把组成原理看成操作系统前面的附录，而能看到二者是一条连续的工程演进线。

\begin{keyidea}
操作系统的大部分核心对象，都可以从硬件限制中推出来：寄存器和 RIP 推出线程上下文，特权级推出系统调用和内核入口，MMU 推出地址空间，定时器推出抢占调度，DMA 和中断推出驱动队列，cache 和多核推出锁、屏障和 per-CPU 数据。
\end{keyidea}

\subsection{第一步：只有一段固件时，还没有操作系统}

从最小机器开始想象：一颗微控制器，一段烧在 Flash 里的固件，一小片 SRAM，几个 MMIO 寄存器控制 GPIO、串口和定时器。复位后 CPU 从固定地址取指，固件设置栈指针，初始化外设，然后进入主循环。这个系统没有进程、没有文件、没有用户态，也没有“内核”这个边界。

\begin{lstlisting}
reset:
  set stack pointer
  initialize clock and UART
  initialize timer
  while true:
    poll input
    update state
    write output
\end{lstlisting}

这种系统能工作，是因为它的信任模型极简单：所有代码由同一个固件作者写入，所有代码拥有同样权限，系统只运行一个应用。若固件写错地址，整机崩溃；若循环太久，外设响应变慢；若中断处理函数和主循环同时改同一个变量，就可能出现竞态。但这些问题还没有被抽象成“操作系统问题”，因为没有多个互不信任程序需要隔离。

此时最重要的硬件概念是复位、时钟、栈、MMIO 和中断。它们已经预告了后来的 OS 机制：复位和启动链会演进成 bootloader 与内核初始化；栈会演进成每线程内核栈；MMIO 会演进成驱动资源；中断会演进成异常入口、调度时钟和设备完成路径。

\subsection{第二步：中断让控制流不再只有一条}

轮询主循环的问题是响应时间受循环长度限制。若主循环正在处理一段耗时计算，串口输入可能溢出，定时事件可能延迟。中断的出现解决了这个问题：硬件事件可以打断当前执行，把 CPU 引到预先登记的入口函数。

\begin{lstlisting}
main loop running:
  instruction N
  timer interrupt becomes pending
  CPU saves minimal interrupted state
  CPU jumps to interrupt vector
  handler runs
  CPU restores state and returns
\end{lstlisting}

这里已经出现了“上下文”的雏形。中断入口必须保存被打断代码需要继续执行的状态，例如 RIP、RFLAGS 和通用寄存器。保存不完整，主程序返回后就会计算错误。保存太多，延迟增加。内核后来的上下文切换，本质上是把这种“保存和恢复执行现场”从中断扩展到多个线程之间。

中断还引入了并发的第一个版本。即使只有一个 CPU，主循环和中断处理函数也会交错访问共享变量。若主循环正在更新一个多字节计数器，中断在中间读取，就可能看到半更新状态。最早的临界区可以简单到关中断：

\begin{lstlisting}
disable_interrupts()
shared_counter = shared_counter + 1
enable_interrupts()
\end{lstlisting}

但这个方法只能保护单核上的中断交错。进入多核以后，另一个 CPU 不受本核关中断影响，必须使用锁和原子操作。理解这一点，才能明白为什么“关中断”和“自旋锁”是不同层次的互斥手段。

\subsection{第三步：定时器把协作式程序推向抢占式系统}

如果所有代码都愿意主动让出 CPU，系统可以用协作式调度：任务运行一会儿，调用 \code{yield}，再让别人运行。这在早期小系统或事件循环里可行，但一旦某个任务死循环，整个系统就失去响应。硬件定时器把 OS 从“相信任务会合作”推进到“内核可以强制收回 CPU”。

\begin{lstlisting}
timer interrupt:
  save current register frame
  account current runtime
  if current consumed its time slice:
    mark current runnable
    choose next runnable task
    switch stack and registers
  return to selected task
\end{lstlisting}

这段路径逼出了三个对象。第一是 task：内核必须有地方保存每个执行流的寄存器、栈、状态和调度账本。第二是 runqueue：内核必须知道哪些 task 现在可运行。第三是 scheduler：当定时器或阻塞事件改变状态时，内核必须选择下一个 task。

定时器同时也定义了时间的可信来源。没有稳定时钟，睡眠超时、网络重传、调度统计、文件时间戳和性能测量都会失真。硬件定时器不是小外设，它是多任务 OS 的节拍器。

\subsection{第四步：特权级把“库函数”变成“内核入口”}

若所有程序都同权，任何程序都可以关中断、改页表、访问磁盘控制器、读写别人的内存。为了运行多个互不信任程序，CPU 必须提供至少两种执行模式：受限的用户态和完全控制硬件的内核态。用户态代码遇到需要特权的操作时，不能直接做，只能通过受控入口请求内核代做。

\begin{lstlisting}
user process wants to read file:
  place syscall number and arguments in registers
  execute syscall/trap instruction
  CPU switches to privileged mode
  CPU jumps to kernel entry
  kernel validates fd, buffer and permission
  kernel performs operation
  CPU returns to user mode with result
\end{lstlisting}

系统调用的本质是特权边界切换，不是普通函数调用。普通函数调用信任调用者传来的指针，系统调用不能信任。用户给内核一个地址，内核必须确认它属于用户地址空间、权限合适、长度不会溢出，并且复制过程中要能处理缺页。若内核直接解引用恶意用户指针，用户程序就可能诱导内核读写任意地址。

这一步还解释了为什么 ABI 很重要。系统调用入口要约定号码放在哪个寄存器，参数放在哪些寄存器，返回值和错误码如何表达，哪些寄存器会被破坏。特权边界不是只靠“进入内核”四个字，它是一套硬件状态、寄存器约定、栈切换、权限检查和返回路径共同组成的协议。

\subsection{第五步：MMU 把地址从数字变成权限对象}

没有 MMU 时，地址就是物理地址。程序写地址 0x1000，就真的写那块物理存储。多个程序共享机器时，这显然不够：一个程序的指针错误可以覆盖另一个程序或内核。MMU 引入虚拟地址到物理地址的翻译，并在翻译时检查权限。

\begin{lstlisting}
CPU executes load/store:
  virtual_address = base + offset
  lookup TLB
  if miss, walk page table
  check present, user/supervisor, read/write/execute
  produce physical address or raise page fault
\end{lstlisting}

这一步把地址空间推出来了。进程不再只是“正在运行的代码”，而拥有一张自己的地址翻译表。相同虚拟地址在不同进程里可以指向不同物理页；内核页可以映射为 supervisor only；代码页可以 executable but not writable；栈页可以 writable but not executable。

MMU 也把异常变成正常控制流。访问不存在的页不一定是错误，可能是懒分配；写只读 COW 页不一定是非法，可能是 fork 后第一次写；访问文件映射缺页可能需要从磁盘读入。page fault handler 是 OS 和硬件合作的典型路径：硬件发现指令无法按当前翻译完成，内核检查 VMA 和 PTE，修复状态后让同一条指令重试。

\subsection{第六步：cache 和 TLB 让“正确顺序”变得昂贵}

cache 和 TLB 是为了速度出现的。没有 cache，CPU 大量时间等待内存；没有 TLB，每次访存都要走多级页表。它们的副作用是状态被复制到了多个地方：内存里有页表，TLB 里有翻译缓存；DRAM 里有数据，多个 CPU cache 里也可能有同一 cache line。

当内核修改页表权限时，只改内存中的 PTE 不够，旧 TLB 项可能还在 CPU 里。当驱动准备 DMA 描述符时，只把字段写到 cache line 不够，设备必须在 doorbell 后能看到正确内容。当两个 CPU 修改同一共享变量时，只看源代码顺序不够，硬件可能通过 store buffer、失效队列和乱序执行改变可见时间。

\begin{lstlisting}
page permission downgrade:
  update PTE in memory
  flush local TLB entry
  send shootdown IPI to other CPUs using this address space
  wait until remote CPUs acknowledge
  only then free or reuse the physical page
\end{lstlisting}

这一步逼出了 TLB shootdown、内存屏障、cache flush、DMA 同步、原子操作和锁。很多 OS 难点并不在“怎样表达一个状态”，而在“怎样让所有硬件观察者在正确时刻看到同一个状态”。

\subsection{第七步：总线和设备把内核推向驱动模型}

只要 CPU、内存和一个串口，系统还能写死设备地址。但真实机器有许多设备，设备可以热插拔，地址窗口由固件或总线枚举分配，中断可以有多个 vector，DMA 地址受设备能力和 IOMMU 限制。内核必须把硬件资源抽象成驱动模型：发现设备、匹配驱动、申请资源、初始化设备、处理中断、提交请求、回收资源。

\begin{lstlisting}
driver probe:
  identify device
  map MMIO BAR or register range
  allocate IRQ/MSI-X vectors
  set DMA mask and create DMA mappings
  initialize submission and completion queues
  register block/net/input/char device object
\end{lstlisting}

驱动模型的关键是生命周期管理，远不止“会写寄存器”。probe 成功后，设备对象对上层可见；remove 时必须阻止新请求、等待或取消旧请求、释放 IRQ、unmap DMA、销毁队列、关闭电源和时钟。若错误路径少释放一步，就泄漏资源；多释放一步，就可能 use-after-free。

\subsection{第八步：DMA 和 completion 把同步调用变成异步协议}

PIO 时代，CPU 可以在一个函数里读完设备数据。DMA 时代，CPU 只提交描述符，设备稍后完成。于是一次 I/O 被拆成两个时间点：提交和完成。中间这段时间里，请求对象、buffer、DMA 映射和等待者都必须保持有效。

\begin{lstlisting}
submit read request:
  pin or allocate buffer pages
  map pages for DMA_FROM_DEVICE
  fill device descriptor
  ring doorbell
  put task to sleep or return EAGAIN

completion interrupt:
  read completion queue
  unmap DMA
  mark buffer uptodate
  wake task or post event
\end{lstlisting}

这个异步协议解释了为什么内核里到处都有 request、bio、skb、urb、kiocb、fence、completion、wait queue。它们都是在保存“已经提交但还没完成”的状态。同步系统调用只是表面：\code{read} 可以在内部提交异步 I/O，然后让当前线程睡眠；设备完成时，中断路径唤醒它，系统调用才返回。

\subsection{第九步：多核把一个内核变成分布式状态机}

单核内核可以用“当前 CPU 正在做什么”理解很多事情。多核内核不行。多个 CPU 同时进入内核，同时修改不同对象，甚至同时修改同一对象。每 CPU runqueue、per-CPU 计数器、本地中断、远程 IPI、RCU grace period、跨 CPU TLB shootdown，都是多核带来的结构。

\begin{longtable}{p{0.22\linewidth}p{0.34\linewidth}p{0.32\linewidth}}
\toprule
多核问题 & 简单做法为什么不够 & OS 机制 \\
\midrule
全局 runqueue 锁竞争 & 所有 CPU 抢同一把锁 & 每 CPU runqueue + 负载均衡 \\
全局引用计数抖动 & cache line 在 CPU 间来回迁移 & per-CPU 计数、批量汇总 \\
读多写少数据 & 普通锁让读者互相阻塞 & RCU、seqlock、读写锁 \\
页表失效传播 & 远端 CPU 仍有旧 TLB & IPI shootdown、PCID \\
NUMA 远端访问 & 内存延迟不均匀 & CPU affinity、页迁移、first touch \\
\bottomrule
\end{longtable}

多核系统里，性能 bug 常常是共享状态设计错误。一个看似无害的全局计数器，在 128 核机器上可能变成所有 CPU 争用的 cache line。一个在单核上正确的“关中断保护”，在多核上不能保护另一个 CPU。OS 设计必须从一开始就问：这个状态属于全局、每 CPU、每进程、每设备，还是每请求？

\subsection{第十步：文件和 socket 是硬件差异的收敛结果}

从硬件看，磁盘、网卡、管道、终端和定时器完全不同。从用户态看，它们可以都变成 fd。这个统一来自 OS 在中间建立的对象层，而不是硬件本身相似：fd 表、open file description、inode/socket/device file、ops 函数表、等待队列和 readiness 事件共同支撑了统一接口。

\begin{lstlisting}
user read(fd):
  fd table maps integer to struct file
  struct file points to object and operations
  operation may hit page cache, pipe buffer, socket receive queue, or device
  if data unavailable:
    task sleeps on object's wait queue
  completion path wakes waiters
\end{lstlisting}

统一接口有巨大价值：shell 重定向、管道、epoll、权限检查、引用计数、close-on-exec 都可以围绕 fd 建立。但统一接口不能抹掉语义差异。普通文件 read 返回 EOF，socket read 受协议影响，pipe 受容量和对端关闭影响，timerfd 返回计数，eventfd 返回计数器。教材若只说“文件是一切”，就会掩盖这些关键差异。

\subsection{从硬件机制到 OS 对象的映射}

下面这张表是第一阶段的压缩地图。后续章节每讲一个 OS 机制，都应能回到这里找到硬件来源。

\begin{longtable}{p{0.22\linewidth}p{0.34\linewidth}p{0.32\linewidth}}
\toprule
硬件机制 & 直接问题 & OS 对象或算法 \\
\midrule
RIP、寄存器、栈 & 执行现场如何保存恢复 & thread context、kernel stack、context switch \\
中断向量 & 异步事件如何进入可信代码 & IDT/vector table、IRQ handler、softirq \\
定时器 & 不合作任务如何被抢占 & scheduler tick、hrtimer、time accounting \\
特权级 & 不可信程序如何请求特权操作 & syscall、trap frame、copy\_from\_user \\
MMU/PTE & 地址如何隔离和授权 & address space、VMA、page fault、COW \\
TLB & 翻译缓存如何保持一致 & shootdown、PCID、lazy TLB \\
cache/coherence & 多核共享数据如何可见 & locks、atomics、barriers、RCU \\
MMIO & CPU 如何控制设备 & ioremap、readl/writel、regmap \\
DMA & 设备如何直接访问内存 & DMA mapping、IOMMU、buffer lifetime \\
MSI/MSI-X & 设备如何报告完成 & IRQ affinity、completion queue、NAPI \\
PCIe/USB/I2C & 设备如何被发现和匹配 & bus model、driver probe、resources \\
DRAM/NUMA & 内存距离和可靠性不同 & page allocator、NUMA policy、memory error handling \\
持久存储 & 写入何时算持久 & page cache、writeback、journal、fsync \\
\bottomrule
\end{longtable}

\subsection{第一阶段读者应该形成的判断力}

到这里，不要求读者已经能写完整内核，但应该能做出一组判断。看到系统调用，要想到特权级、寄存器 ABI、用户指针检查和返回路径。看到 page fault，要想到 MMU、VMA、PTE、TLB、COW 和缺页重试。看到设备驱动，要想到 MMIO、DMA、IRQ、队列、屏障和错误路径。看到多核性能问题，要想到 cache line、锁粒度、per-CPU、IPI 和 NUMA。

这就是第一阶段的目标：不把硬件当背景，不把 OS 当概念，而是把二者合成一条可追踪的执行链。
